\subsection{FFT}

Seja $A(x) = a_0 + a_1 x + \dots + a_{n-1} x^{n-1}$ e $B(x)$ definido analogamente. Queremos saber os coeficientes de $A(x) \cdot B(x)$ em $O(n \log n)$.

Para isso, $\text{DFT(A)} = \text{DFT}(a_0, \dots, a_{n-1}) =  (y_0, \dots, y_{n-1}) = (A(w_n^0), \dots, A_(w_n^{n-1}))$. Ademais, defina $\text{InverseDFT}(y_0, \dots, y_{n-1}) = (a_0, \dots, a_{n-1})$. Note que $AB = \text{InverseDFT}(\text{DFT}(AB)) = \text{InverseDFT}(\text{DFT}(A)\text{DFT}(B))$ (em que a multiplicação de dentro é component-wise). Sendo assim, a ideia é fazer a DFT e a InverseDFT em $O(n \log n)$, para daí obter $AB$.

Essa FFT não é nada boa para multiplicações envolvendo coeficientes altos. Vão ter muitos erros de imprecisão. Geralmente se não tiver gente perto de $10^{15}$ (pra double) ou $10^{18}$ (pra long double) ta de boa.

\begin{lstlisting}[language=c++, breaklines=true]
using cd = complex<double>;
const double PI = acos(-1);

void fft(vector<cd> & a, bool invert) {
    int n = a.size();

    for (int i = 1, j = 0; i < n; i++) {
        int bit = n >> 1;
        for (; j & bit; bit >>= 1)
            j ^= bit;
        j ^= bit;

        if (i < j)
            swap(a[i], a[j]);
    }

    for (int len = 2; len <= n; len <<= 1) {
        double ang = 2 * PI / len * (invert ? -1 : 1);
        cd wlen(cos(ang), sin(ang));
        for (int i = 0; i < n; i += len) {
            cd w(1);
            for (int j = 0; j < len / 2; j++) {
                cd u = a[i+j], v = a[i+j+len/2] * w;
                a[i+j] = u + v;
                a[i+j+len/2] = u - v;
                w *= wlen;
            }
        }
    }

    if (invert) {
        for (cd & x : a)
            x /= n;
    }
}

vector<int> multiply(vector<int> const& a, vector<int> const& b) {
    vector<cd> fa(a.begin(), a.end()), fb(b.begin(), b.end());
    int n = 1;
    while (n < a.size() + b.size()) 
        n <<= 1;
    fa.resize(n);
    fb.resize(n);

    fft(fa, false);
    fft(fb, false);
    for (int i = 0; i < n; i++)
        fa[i] *= fb[i];
    fft(fa, true);

    vector<int> result(n);
    for (int i = 0; i < n; i++)
        result[i] = round(fa[i].real());
    return result;
}
\end{lstlisting}

\subsubsection{Number theoretic transform}

\textbf{Objetivo:} multiplicar polinômios agora com os coeficientes módulo $p$. 

\textbf{Solução:} para $p=73400233$, $w_{2^{20}}=5$ é uma $2^{20}$-ésima raiz primitiva. Podemos, então, utilizar a seguinte FFT:

\begin{lstlisting}[language=c++, breaklines=true]
/**
 * NTT eh o mesmo que FFT, so que com um *n-th root* diferente
 * MOD tem que ser um primo do tipo p = c * 2^k + 1
 * Com isso, a n-th root existe pra n = 2^k. 
 * Essa n-th root é g^c, onde g é um *primitive root* de 2^k
 */

namespace NTT {
const int mod = 	998244353;// 998244353	7340033
const int root = 	363395222;// 15311432	5
const int root_1 = 	704923114;// 469870224	4404020
const int root_pw = 1 << 19;//	 1 << 23;	1 << 20;
//					~5*10^5	 	 ~8*10^6	~10^6

int fast_pow(int a, int b, int m) {
    int ans = 1;
    while(b) {
        if(b&1) ans = 1LL * ans * a % m;
        a = 1LL * a * a % m;
        b >>= 1;
    }
    return ans;
}
int inverse(int x, int m) { return fast_pow(x, m - 2, m); }

/* INICIO */
/** So quando as constantes acima nao forem suficientes */

map<int, int> factor(int n) {
    map<int, int> ans;
    for(int i = 2; i * i <= n; ++i) {
        while(n%i == 0) {
            ans[i]++;
            n /= i;
        }
    }
    if(n > 1) ans[n]++;
    return ans;
}

int Phi(int n) {
    auto fact = factor(n);
    int ans = 1;

    for(auto [a, b] : fact) {
        ans *= fast_pow(a, b, n + 1) - fast_pow(a, b-1, n + 1);
    }
    return ans;
}

int prim_root(int n) {
    int phi = Phi(n);
    auto fact = factor(phi);

    for (int res=2; res<=n; ++res) {
        if(__gcd(res, n) > 1) continue;
        bool ok = true;
        for (auto [a, b] : fact) {
            ok &= fast_pow(res, phi / a, n) != 1;
            if(!ok) break;
        }
        if (ok) return res;
    }
    return -1;
}

// Generates NTT constants for any Mod and prints it on the screen
void generate(int Mod) {

    int n = 1;
    int aux = Mod-1;
    while((aux&1) == 0) aux >>= 1, n <<= 1;
    int c = aux;

    // g = primitive root de Mod
    int g = prim_root(Mod);
    if(g == -1) {
        printf("No constants could be found, cant find primitive root of %d\n", n);
        return;
    }
    int root = fast_pow(g, c, Mod);
    int root_1 = inverse(root, Mod);
    printf("mod = %d, root = %d, root_1 = %d root_pw = %d\n", Mod, root, root_1, n);
    assert(fast_pow(root, n, Mod) == 1);
    // g^c
}
/* FIM */


inline void ntt(vector<int>& a, bool invert) {
    int n = a.size();
    
    for (int i = 1, j = 0; i < n; i++) {
        int bit = n >> 1;
        for (; j & bit; bit >>= 1)
            j ^= bit;
        j ^= bit;

        if (i < j)
            swap(a[i], a[j]);
    }

    for (int len = 2; len <= n; len <<= 1) {
        int wlen = invert ? root_1 : root;

        for (int i = len; i < root_pw; i <<= 1)
            wlen = 1LL * wlen * wlen % mod;

        for (int i = 0; i < n; i += len) {
            int w = 1;
            for (int j = 0; j < len / 2; j++) {
                int u = a[i+j], v = 1LL * a[i+j+len/2] * w % mod;
                a[i+j] = u + v < mod ? u + v : u + v - mod;
                a[i+j+len/2] = u - v >= 0 ? u - v : u - v + mod;
                w = 1LL * w * wlen % mod;
            }
        }
    }


    if (invert) {
        int n_1 = inverse(n, mod);
        for (int & x : a)
            x = 1LL * x * n_1 % mod;
    }
}

vector<int> multiply(vector<int> const& a, vector<int> const& b) {
    vector<int> ca(a.begin(), a.end()), cb(b.begin(), b.end());
    int n = 1;
    while(n < (int) max(a.size(), b.size())) n <<= 1;
    n <<= 1;

    ca.resize(n); cb.resize(n);
    
    ntt(ca, false); ntt(cb, false);
    
    for (int i = 0; i < n; i++) ca[i] = 1LL * ca[i] * cb[i] % mod;
    ntt(ca, true);

    return ca;
}
};
\end{lstlisting}

Para módulos da forma $p = c 2^k + 1$ (com $p$ primo), $g^c$ é uma $2^k$-ésima raiz da unidade, em que $g$ é uma raiz primitiva módulo $p$.

\textbf{Atenção:} Se em algum momento ela tentar calcular um fft de um vetor de tamanho maior que $2^k$, vai bugar!

Como calcular NTT com módulo $M$ arbitrário? Uma opção é usar teorema chinês dos restos em primos da forma $c2^k+1$

Uma outra opção é decompor os polinômios $A(x)$ e $B(x)$ em $A = A_1 + A_2 \cdot C$ e $B = B_1 + B_2 \cdot C$ em que $C \approx \sqrt{M}$ (de modo que os coeficientes de cada um deles sejam pequenos) e daí
$$ A \cdot B = A_1 \cdot B_1 + (A_1 \cdot B_2 + A_2 \cdot B_1) \cdot C + (A_2 \cdot B_2) \cdot C^2$$
ou seja, com isso, os coeficientes calculados (de cada produto $A_i B_j$) não passam de $n \cdot M$.

\prob{1}{fft1} Número de pares de números com uma dada soma.

\textbf{Solução:} Eleve ao quadrado o polinômio de frequências.

\prob{2}{fft2} Produto interno entre array a e shifts cíclicos de array b

\textbf{Solução:} Reverte a, completa com zeros, apenda b em b e calcula o produto entre os novos a e b. Na array c, c[k] = produto escalar entre a e o left shift cíclico k-(n-1) de b (com k indo de n-1 a 2n-2).

\begin{lstlisting}[language=c++, breaklines=true]
// NAO E exatamente esse problema
// recebo duas fitas de DNA |s| >= |r|
// determinar r.size() - maior produto interno
// ou seja eu vou deslizando r em s (sem passar dos limites)
#include <bits/stdc++.h>

using namespace std;

using cd = complex<double>;
const double PI = acos(-1);
typedef long long int lli;

void fft(vector<cd>& a, bool invert) {}

vector<lli> multiply(vector<lli> const& a, vector<lli> const& b) {}

int conversion_map[200];

int main() {
    string s;
    string r;
    cin >> s;
    cin >> r;
    conversion_map['A'] = 0;
    conversion_map['C'] = 1;
    conversion_map['T'] = 2;
    conversion_map['G'] = 3;

    vector<lli> pols_s[4];
    vector<lli> pols_r[4];
    vector<lli> prods[4];
    for (int i = 0; i < s.size(); i++) {
        int atual = conversion_map[s[i]];
        for (int j = 0; j < 4; j++) {
            pols_s[j].push_back(atual == j);
        }
    }

    for (int i = 0; i < r.size(); i++) {
        int atual = conversion_map[r[i]];
        for (int j = 0; j < 4; j++) {
            pols_r[j].push_back(atual == j);
        }
    }

    for (int j = 0; j < 4; j++) {
        reverse(pols_r[j].begin(), pols_r[j].end());
        prods[j] = multiply(pols_s[j], pols_r[j]);
        if (prods[j].size() < s.size() - 1) prods[j].resize(s.size() - 1);
    }

    int max_ans = -1;

    for (int i = r.size() - 1; i <= s.size() - 1; i++) {
        int curr_ans = 0;
        for (int j = 0; j < 4; j++) {
            curr_ans += prods[j][i];
        }
        if (curr_ans > max_ans) {
            max_ans = curr_ans;
        } else if (curr_ans < max_ans)
            continue;
    }

    cout << r.size() - max_ans << "\n";
}
\end{lstlisting}

\prob{3}{fft3} Temos duas fitas de booleanos a e b. Queremos encontrar o número de maneiras de aplicar uma na outra de modo que nenhum par de uns estejam na mesma posição.

\textbf{Solução:} análoga a \ref{prob:fft2}.

\prob{4}{fft4} Dado um texto $T$, um padrão $P$ (os 2 de letra minuscula), calcular todas as ocorrências do padrão no texto.

\textbf{Solução:} Criamos
$$A(x) = a_0 + \dots + a_{n-1} x^{n-1}, \, n= |T|$$
com $a_i = e^{i \alpha_i}$ onde $\alpha_i = \dfrac{2 \pi T[i]}{26}$
$$B(x) = b_0 + \dots + b_{m-1} x^{m-1}, \, m = |P|$$
com $b_i = e^{-i \beta_i}$ onde $\beta_i = \dfrac{2 \pi P[m-i-1]}{26}$

Agora, considere $C(x) = A(x) \cdot B(x)$. Se $c_{m-1-i} = m$, é porque teve um match a partir de $T[i]$.

\prob{5}{fft5} Permitindo wildcards no padrão (letras que dão match com qualquer uma)

\textbf{Solução:} basta  setar $b_i=0$ sempre que $P[m-i-1] = *$. Aí se o padrão tiver $x$ wildcards, então $c_{m-1-i} = m-x$ indica que a partir do índice $T[i]$ coincidem.

\subsubsection{FFT 2D}

\begin{lstlisting}[language=c++, breaklines=true]
const int B = 12; // nesse problema os pols envolvidos tinham no max x^800 y^800
const int N = 1 << B;

cd roots[N];
int inv[N];

cd get_root(int k, int n) { return roots[k * (N / n)]; }

void precalc() {
    for (int i = 0; i < N; i++) {
        double ang = 2 * PI * i / N;
        roots[i] = cd(cos(ang), sin(ang));
    }
}

bool bit(int mask, int i) { return (mask >> i) & 1; }

void precalc_inv(int n) {
    int b = 0;
    while ((1 << b) < n)
        ++b;
    assert((1 << b) == n);
    assert(b <= B);

    inv[0] = 0;
    int hb = -1;
    for (int i = 1; i < n; ++i) {
        if (bit(i, hb + 1)) {
            ++hb;
        }
        inv[i] = inv[i ^ (1 << hb)] ^ (1 << (b - hb - 1));
    }
}

// ****copia e cola fft aq

void fft_2d(vector<vector<cd>>& a, bool rev) {
    precalc_inv((a[0].size()));
    for (auto& row : a) {
        fft(row, rev);
    }
    precalc_inv((a.size()));
    for (int j = 0; j < a.front().size(); j++) {
        vector<cd> col;
        for (int i = 0; i < a.size(); i++)
            col.push_back(a[i][j]);
        fft(col, rev);
        for (int i = 0; i < a.size(); i++)
            a[i][j] = col[i];
    }
}

vector<vector<cd>> mult(vector<vector<cd>> x, vector<vector<cd>> y) {
    int b_rows = 0;
    while ((1 << b_rows) <= max(x.size(), y.size()))
        ++b_rows;
    ++b_rows;

    int b_cols = 0;
    while ((1 << b_cols) <= max(x.front().size(), y.front().size()))
        ++b_cols;
    ++b_cols;

    x.resize(1 << b_rows);
    y.resize(1 << b_rows);
    for (auto& row : x) {
        row.resize(1 << b_cols, 0);
    }
    for (auto& row : y) {
        row.resize(1 << b_cols, 0);
    }

    fft_2d(x, 0);
    fft_2d(y, 0);

    for (int i = 0; i < x.size(); i++)
        for (int j = 0; j < x[i].size(); j++)
            x[i][j] *= y[i][j];

    fft_2d(x, 1);

    return x;
}
\end{lstlisting}
